\documentclass{article}

\usepackage{iclr2027_conference,times}
\usepackage{hyperref}
\usepackage{url}
\usepackage{graphicx}
\usepackage{booktabs}
\usepackage{amsmath}
\usepackage{amssymb}
\usepackage{array}
\usepackage{multirow}
\usepackage{xcolor}
\usepackage{etoolbox}

\newcommand{\method}{\textsc{SEEK}}
\newcommand{\tailleak}{\textsc{Tail Leakage}}
\newcommand{\tl}{\mathrm{TL}}
\newcommand{\hotpot}{HotpotQA}
\newcommand{\twowiki}{2Wiki}
\newcommand{\musique}{MuSiQue}
\newcommand{\popqa}{PopQA}
\newcommand{\emf}[2]{#1/#2}


\title{Tail Leakage: A Simple Confidence Signal for Adaptive Retrieval}

\author{Anonymous authors\\
Working draft --- not submitted or published}

% The official anonymous style hard-codes "Paper under double-blind review"
% instead of using \author. Keep the style file pristine and replace that
% status locally until this manuscript is actually submitted.
\makeatletter
\patchcmd{\@maketitle}
  {Anonymous authors\\Paper under double-blind review}
  {Anonymous authors\\Working draft --- not submitted or published}
  {}{\PackageError{seek-draft}{Could not patch draft status}{}
  }
\makeatother

% Keep this commented for the anonymous submission.
% \iclrfinalcopy

\begin{document}

\maketitle

% Repository-status banner: remove only if this manuscript is actually
% submitted. The official review style otherwise labels drafts "under review."
\lhead{Anonymous ICLR 2027 draft --- not submitted or published}
\begin{center}
  \textbf{Anonymous draft --- not submitted or published (2026-08-24)}
\end{center}

\begin{abstract}
Adaptive retrieval requires a model to decide when its internal knowledge is
insufficient. Token entropy is a natural training-free signal, but it mixes
competition among plausible leading tokens with probability mass dispersed
across the rest of the vocabulary. We introduce \tailleak{}, a confidence score
that retains only the latter effect. At each generated position, it measures
the negative log probability mass retained by the model's top-\(k\) tokens and
sums this quantity over the answer. The score is available from one greedy
decoding pass and requires neither sampling nor a learned router. We study
\tailleak{} in \method{}, a dual-path open-domain question-answering framework
that applies the score to sub-question retrieval and preserves both direct and
decomposition candidates. Across four development-set traces with
Llama-3-8B-Instruct, \tailleak{} reaches \(0.749\) macro AUROC for detecting
wrong direct answers, compared with \(0.726\) for maximum token entropy. In the
tested end-to-end configuration, replacing entropy with \tailleak{} raises
macro EM/F1 from \(40.52/47.61\) to \(41.88/48.90\). Decomposition-path queries
also increase answer-containing retrieval recall on three multi-hop datasets,
while leaving single-hop PopQA nearly unchanged. These results support
\tailleak{} as a useful ranking and gating signal, but confidently wrong
answers remain a fundamental limitation.
\end{abstract}

\section{Introduction}
\label{sec:introduction}

Retrieval-augmented generation (RAG) gives a language model access to
information beyond its parameters by conditioning generation on external
passages~\citep{rag_survey}. A fixed retrieval policy is nevertheless poorly
matched to heterogeneous questions. Retrieval can add latency and distracting
evidence when the model already knows a simple fact, yet a single query is
often insufficient for a multi-hop question whose decisive entity becomes
known only after an intermediate step~\citep{ircot,ren2023}. Adaptive RAG
therefore needs to decide not only \emph{what} to retrieve, but also
\emph{when} the current reasoning state warrants retrieval.

One appealing solution is to use the language model's own token distribution.
Uncertainty-triggered methods can act during generation without a separately
trained routing model~\citep{flare,dragin,targ}. In practice, however,
uncertainty is not synonymous with factual error. A model can distribute
probability among several harmless aliases or formatting variants while
remaining factually correct; it can also assign almost all probability to a
fluent hallucination. Standard entropy combines these cases because it records
both competition among high-probability candidates and dispersion into the
large low-probability tail.

This paper asks whether a more selective statistic of the token distribution
can provide a better confidence signal for adaptive retrieval. We introduce
\tailleak{}, which measures how much probability escapes the model's
own small set of leading candidates. If \(m_t\) is the probability mass of the
top-\(k\) tokens at generation step \(t\), the answer-level score is
\[
    \tl(a)=-\sum_{t\in a}\log m_t.
\]
The score is near zero when a small candidate set retains almost all
probability and increases monotonically as mass moves into the tail. Unlike
full entropy, it discards how probability is divided inside the leading set
and how the remaining mass is arranged within the tail. It is computed from
the logits already produced during greedy decoding, without extra generations,
labels, or router training.

We evaluate this signal inside \method{}, short for Self-certainty Evaluated
Evidence-gated Knowledge-retrieval. This dual-path framework for open-domain
question answering first
produces a direct answer and also decomposes the original question into
ordered, executable sub-questions. Each sub-question is tentatively answered
without retrieval; \tailleak{} then decides whether to retrieve and re-answer
that step. The decomposed steps are synthesized into a second candidate, and a
transparent late rule selects between the direct and decomposition candidates.
This design makes decomposition reversible: it can expose bridge entities for
retrieval without discarding the original-question answer.

Our evaluation uses four development-set traces spanning three multi-hop
datasets---\hotpot{}, \twowiki{}, and \musique{}---and the single-hop
\popqa{} benchmark. Among six representative confidence families,
\tailleak{} gives the highest macro AUROC for detecting wrong direct answers
(\(0.749\), versus \(0.726\) for maximum token entropy). In the tested
end-to-end configuration, the \tailleak{} system obtains \(41.88/48.90\)
macro EM/F1, compared with \(40.52/47.61\) when maximum token entropy is used.
The accuracy-first retrieval gate is effectively tied with always retrieving
while skipping \(6.16\%\) of decomposed steps; more aggressive thresholds show
a gradual accuracy--retrieval trade-off. The results do not imply that the
score is a correctness test: confidently wrong answers remain, and final
routing contributes only a small net gain.

Our contributions are:
\begin{itemize}
    \item We derive \tailleak{}, a simple, monotonic, training-free statistic
    that isolates probability mass leaving the leading token set and
    accumulates it across an answer.
    \item We instantiate the score as a sub-question retrieval gate in
    \method{}, a reversible dual-path framework for decomposition-based RAG.
    \item We compare representative confidence signals, retrieval behavior,
    and end-to-end QA across four datasets, while identifying the boundary
    created by confident errors and imperfect path routing.
\end{itemize}

\section{Related Work}
\label{sec:related}

\paragraph{Adaptive retrieval.}
Adaptive-RAG routes questions to retrieval strategies using a learned
complexity classifier~\citep{adaptiverag}, while SKR estimates whether a model
already knows an answer before retrieval~\citep{skr}. FLARE instead drafts the
next sentence and retrieves when low-confidence tokens appear~\citep{flare};
DRAGIN combines a generation-time trigger with attention-based query
formulation~\citep{dragin}. TARG studies a lightweight question-level decision
from prefix logits, including entropy and top-logit margins~\citep{targ}.
These methods make uncertainty useful operationally, but most apply it either
before reasoning or at the free-form token/sentence level. \method{} applies a
training-free signal at explicit sub-question boundaries.

Learned policies provide finer control. Self-RAG trains reflection tokens that
control retrieval and evaluate evidence~\citep{selfrag}. DeepRAG explores
retrieve-or-answer decisions at individual sub-queries and calibrates the
policy from preference data~\citep{deeprag}. Such policies can jointly learn
reasoning and retrieval behavior, whereas our objective is complementary: a
score that can be inspected and attached to an existing greedy generation
without policy training.

\paragraph{Decomposition and multi-step retrieval.}
IRCoT interleaves chain-of-thought sentences with retrieval~\citep{ircot}.
Question-decomposition methods instead expose explicit intermediate
operations or sub-questions~\citep{decomprc,break_qdmr}. Decomposition can
create sharper retrieval queries, but a wrong bridge answer can corrupt later
queries and short sub-questions can omit constraints from the original
question. \method{} therefore retains a direct candidate and keeps
original-question evidence available to the decomposition path. We treat these
components as an evaluation framework for the confidence signal rather than
attributing all system gains to decomposition alone.

\paragraph{Confidence and uncertainty in language models.}
Sequence probability, token entropy, verbalized self-evaluation such as
P(True), and disagreement across sampled answers capture different aspects of
model confidence~\citep{kadavath2022,semanticuncertainty,semanticentropy}.
They are useful ranking signals, but a model may remain confidently wrong.
\tailleak{} is closest to logit-based uncertainty methods, yet differs in what
it retains: it ignores the exact probability arrangement within the head and
tail and records only the mass crossing their boundary. This choice follows
from an entropy decomposition in Section~\ref{sec:tail-leakage} and is tested
against representative alternatives in Section~\ref{sec:confidence-results}.

\section{Tail Leakage}
\label{sec:tail-leakage}

Let \(p_t(v)\) be the next-token distribution at position \(t\) during
greedy generation, and let \(V_k(p_t)\) contain its \(k\) most probable
tokens. We define the retained mass
\begin{equation}
    m_t^{(k)}=\sum_{v\in V_k(p_t)}p_t(v),
    \qquad
    \ell_t^{(k)}=1-m_t^{(k)},
    \label{eq:retained-mass}
\end{equation}
where \(\ell_t^{(k)}\) is the mass that leaks outside the leading set. We use
\(k=5\) throughout the main experiments and write \(m_t\) for
\(m_t^{(5)}\).

\subsection{Why isolate leaked mass?}

Token entropy measures the total spread of \(p_t\), but this spread has
several sources. Let \(\hat p_t^{\,\mathrm{in}}\) and
\(\hat p_t^{\,\mathrm{out}}\) denote the distributions obtained by
renormalizing the tokens inside and outside \(V_5(p_t)\), respectively. The
entropy chain rule gives
\begin{align}
H(p_t)={}&-m_t\log m_t-(1-m_t)\log(1-m_t) \nonumber\\
         &+m_t H(\hat p_t^{\,\mathrm{in}})
          +(1-m_t)H(\hat p_t^{\,\mathrm{out}}).
\label{eq:entropy-decomposition}
\end{align}
The first term records how much mass crosses the top-five boundary. The other
terms record competition among the leading candidates and the arrangement of
probability within the tail. For short-answer QA, competition among leading
tokens can reflect casing, aliases, or tokenization rather than factual
uncertainty. Conversely, the internal shape of a very large tail may change
without materially changing the model's leading hypothesis. Equation
\ref{eq:entropy-decomposition} motivates testing the boundary-crossing mass
separately instead of compressing all three effects into entropy.

\subsection{Token- and answer-level score}

We convert retained mass into a monotonic risk score:
\begin{equation}
    \tl_t=-\log m_t,
    \qquad
    \tl(a)=\sum_{t\in\mathcal{T}(a)}\tl_t
          =-\log\prod_{t\in\mathcal{T}(a)}m_t ,
\label{eq:tail-leakage}
\end{equation}
where \(\mathcal{T}(a)\) contains the generated answer tokens and excludes
EOS. The negative logarithm has two useful properties. First, it increases
strictly as mass leaves the top-five set; unlike binary entropy, it does not
decrease once leaked mass exceeds \(0.5\). Second, it is additive across
positions, allowing small amounts of leakage distributed over a multi-token
answer to accumulate.

\tailleak{} requires only the top-\(k\) probabilities already available during
a greedy pass. It is therefore cheaper than sample-disagreement or semantic
entropy estimates that require multiple generations. It is a ranking score,
not a calibrated probability that an answer is wrong. In particular, if a
model concentrates almost all probability on a false continuation,
Equation~\ref{eq:tail-leakage} provides no warning.

\section{Adaptive Retrieval with \method{}}
\label{sec:seek}

We use \method{} to test whether \tailleak{} is useful beyond offline error
detection. Given a question \(q\), a language model \(\mathcal{M}\), and a
passage corpus \(\mathcal{P}\), \method{} preserves two answer candidates until
late routing:

\begin{figure}[t]
\centering
\small
\setlength{\fboxsep}{6pt}
\fbox{\parbox{0.93\linewidth}{
\textbf{Direct path:}
\(q \rightarrow\) retrieve with \(q\) \(\rightarrow a^{\mathrm{dir}}\)
\(\quad\)\(\Vert\)\(\quad\)
\textbf{Decomposition path:}
\(q\rightarrow(s_1,\ldots,s_n)\rightarrow\) gate retrieval at each \(s_i\)
\(\rightarrow a^{\mathrm{dec}}\).
\par\medskip
\centering
\(a^{\mathrm{dir}},a^{\mathrm{dec}}\rightarrow\)
transparent late routing \(\rightarrow\hat a\)
}}
\caption{The two candidate paths in \method{}. Decomposition is retained as a
retrieval strategy rather than an irreversible early decision.}
\label{fig:seek-overview}
\end{figure}

\paragraph{Executable decomposition.}
A LoRA-adapted decomposer maps \(q\) to ordered, knowledge-graph-style
sub-questions. A step has the form
\(\texttt{entity >> relation}\); a placeholder such as \(\texttt{\#1}\)
refers to the answer of an earlier step. For example,
\(\texttt{Inception >> director}\) followed by
\(\texttt{\#1 >> place of birth}\) becomes a directly executable second query
after the first answer is substituted. One-step reformulations are allowed,
which prevents the representation from forcing a multi-hop chain onto
single-hop questions.

\paragraph{Step-level retrieval gate.}
For each resolved sub-question \(s_i\), the model first generates a tentative
answer \(\tilde a_i\) without retrieved evidence and computes
\(\tl(\tilde a_i)\). Retrieval is invoked when
\begin{equation}
    \pi(s_i)=\mathbb{I}\!\left[\tl(\tilde a_i)\geq\tau\right].
\label{eq:gate}
\end{equation}
If the gate fires, BM25 returns passages and the model re-answers \(s_i\);
otherwise \(\tilde a_i\) is retained. Thus the signal operates at a natural
reasoning boundary rather than at an arbitrary sentence or token position.

\paragraph{Context preservation and synthesis.}
Short entity-centered sub-questions are useful retrieval queries but can omit
comparison targets or disambiguating constraints. When a step retrieves,
\method{} therefore deduplicates the union of passages retrieved from \(s_i\)
and from the original question \(q\). After all placeholders are resolved, a
base-model synthesis pass reads \(q\) and the ordered \((s_i,a_i)\) pairs to
produce \(a^{\mathrm{dec}}\). The direct candidate \(a^{\mathrm{dir}}\)
answers \(q\) from original-question evidence.

\paragraph{Late routing.}
A deterministic verifier defaults to the decomposition candidate and switches
to \(a^{\mathrm{dir}}\) only when direct-answer \tailleak{} is low and the
decomposition path shows instability, such as high post-synthesis leakage,
malformed placeholders, or an incompatible answer type. The verifier does not
regenerate an answer or train a policy. We include it to preserve complementary
candidates, but do not treat it as the principal source of accuracy:
Section~\ref{sec:routing-results} shows that its net effect is small. Exact
rules and diagnostic counts appear in Appendix~\ref{app:details}.

\section{Experimental Setup}
\label{sec:experiments}

\paragraph{Datasets.}
We evaluate on the first \(1{,}000\) development questions, without shuffling,
from each of \hotpot{}~\citep{hotpotqa}, \twowiki{}~\citep{twowiki},
\musique{}~\citep{musique}, and \popqa{}~\citep{popqa}. The first three
benchmarks contain multi-hop questions; \popqa{} consists of single-hop,
entity-centric factual questions. The chosen \musique{} slice contains only
two-hop questions because the development set is ordered by hop count. We
therefore report these results as a controlled development-set study, not as
full-benchmark estimates.

\paragraph{Retrieval and generation.}
The corpus is an Elasticsearch BM25 index of approximately \(21\) million
non-overlapping 100-word English Wikipedia passages following the DPR corpus
construction~\citep{dpr}. Retrieval returns eight passages. Original questions
match paragraph text; short decomposition queries match title and paragraph
text. Each retrieved decomposition step also receives deduplicated evidence
from the original question. This path-aware configuration is fixed across
datasets. Because original and decomposed queries use different title fields,
the retrieval comparison in Section~\ref{sec:retrieval-results} measures the
complete decomposition-path query policy rather than decomposition wording in
isolation.

All components use Llama-3-8B-Instruct with greedy decoding. The decomposer is
trained once on \(12{,}000\) examples: \(4{,}000\) each from \musique{} gold
sub-questions, \twowiki{} evidence-triple conversions, and \popqa{} one-step
triples. Evaluation IDs are excluded from the \popqa{} training portion.
Training uses LoRA~\citep{lora} with rank \(16\), scale \(32\), dropout \(0.05\),
and all linear modules for three epochs. We use bf16, a 1,024-token context,
effective batch size \(16\), learning rate \(10^{-4}\), cosine decay, and
\(10\%\) warmup in LLaMA-Factory~\citep{llamafactory}. The adapter is disabled
during final synthesis.

\paragraph{Operating point and metrics.}
The main retrieval gate uses \(k=5\) and the accuracy-first threshold
\(\tau=0.002305\), which retrieves \(93.84\%\) of decomposed steps. This point
was selected from a pooled threshold sweep on the same development traces.
Accordingly, the threshold sweep is descriptive and should not be interpreted
as held-out hyperparameter selection; disjoint calibration is a required
confirmatory experiment. We report normalized exact match (EM), token-level
F1, pooled step retrieval rate, answer-containing Hit@5, and AUROC for
predicting whether a direct answer is wrong. Generation is deterministic and
we do not report seed variance.

\paragraph{Comparisons.}
For confidence analysis, we compare \tailleak{} with maximum token entropy,
sequence negative log-likelihood, sampled-answer disagreement, the TARG
top-logit margin~\citep{targ}, and P(True) self-evaluation~\citep{kadavath2022}.
For QA, we include Llama-3-8B results reported by DeepRAG~\citep{deeprag} for
CoT, CoT-Retrieve, FLARE, DRAGIN, and DeepRAG. Those values were not rerun in
our harness, so implementation, prompting, and normalization differences
prevent a strict controlled ranking. All \method{} ablations use saved traces
from our own pipeline.

\section{Results}
\label{sec:results}

\subsection{End-to-end question answering}
\label{sec:qa-results}

Table~\ref{tab:main-results} compares \method{} with published
Llama-3-8B results. \method{} has the highest reported average in this set and
the strongest values on \twowiki{}, \musique{}, and \popqa{}, while remaining
below DeepRAG on \hotpot{}. The comparison is encouraging but not controlled:
the baseline rows come from the DeepRAG paper and may differ in prompts,
implementation, and scoring. We therefore do not claim a new state of the art.
On \popqa{}, the decomposition candidate is a one-step query reformulation,
not a multi-hop reasoning chain.

\begin{table}[t]
\caption{EM/F1 compared with reported Llama-3-8B results.}
\label{tab:main-results}
\centering
\scriptsize
\setlength{\tabcolsep}{2.6pt}
\resizebox{\linewidth}{!}{%
\begin{tabular}{lccccc}
\toprule
Method & \hotpot{} & \twowiki{} & \musique{} & \popqa{} & Macro \\
\midrule
CoT\(^\dagger\) & 27.20/37.75 & 28.20/34.85 & 13.70/22.97 & 21.20/25.33 & 22.58/30.23 \\
CoT-Retrieve\(^\dagger\) & 34.90/46.85 & 35.80/43.41 & 19.10/28.18 & 32.80/45.87 & 30.65/41.08 \\
FLARE\(^\dagger\) & 23.80/32.88 & 30.30/37.45 & 14.50/23.57 & 28.80/40.61 & 24.35/33.63 \\
DRAGIN\(^\dagger\) & 27.60/38.05 & 29.10/35.68 & 11.80/19.97 & 22.60/28.53 & 22.78/30.56 \\
DeepRAG\(^\dagger\) & \textbf{40.70}/\textbf{51.54} & 48.10/53.25 & 22.50/30.40 & 42.50/47.80 & 38.45/45.75 \\
\method{} & 37.80/49.03 & \textbf{49.30}/\textbf{56.22} & \textbf{26.50}/\textbf{34.63} & \textbf{53.90}/\textbf{55.73} & \textbf{41.88}/\textbf{48.90} \\
\bottomrule
\end{tabular}}
\vspace{2pt}
\parbox{\linewidth}{\scriptsize \(\dagger\) Values reported by
\citet{deeprag}; not rerun in our harness.}
\end{table}

\subsection{Does Tail Leakage identify answer risk?}
\label{sec:confidence-results}

Table~\ref{tab:confidence} reports AUROC for predicting a wrong direct answer;
larger values indicate better ranking. \tailleak{} has the strongest macro
AUROC among the tested representatives (\(0.749\)), improving over maximum
token entropy and sequence NLL (both \(0.726\)). Its advantage is not uniform
across every dataset: sequence NLL ties it on \popqa{}, and P(True) ties it on
\musique{}. The correct conclusion is therefore limited to these traces and
the tested signal families.

\begin{table}[t]
\caption{AUROC for detecting wrong direct answers.}
\label{tab:confidence}
\centering
\scriptsize
\setlength{\tabcolsep}{3.6pt}
\begin{tabular}{lccccc}
\toprule
Signal & \hotpot{} & \twowiki{} & \musique{} & \popqa{} & Macro \\
\midrule
\tailleak{} & \textbf{.757} & \textbf{.717} & \textbf{.726} & \textbf{.798} & \textbf{.749} \\
Max token entropy & .743 & .671 & .711 & .780 & .726 \\
Sequence NLL & .733 & .667 & .707 & \textbf{.798} & .726 \\
Sample disagreement & .678 & .659 & .680 & .674 & .672 \\
TARG margin & .669 & .589 & .618 & .682 & .639 \\
P(True) & .655 & .567 & \textbf{.726} & .579 & .632 \\
\bottomrule
\end{tabular}
\end{table}

The end-to-end comparison points in the same direction. Using maximum token
entropy gives \(40.52/47.61\) macro EM/F1, while the \tailleak{}
configuration gives \(41.88/48.90\). The gain is again non-uniform:
maximum entropy is \(0.70\) EM higher on \twowiki{}, while \tailleak{} is
higher on the other three datasets. An entropy-component analysis in
Appendix~\ref{app:confidence} provides mechanism-level evidence: the
head-versus-tail mass component reaches \(0.742\) macro AUROC, whereas the
arrangement within the tail is near chance at \(0.502\).

\subsection{Retrieval behavior}
\label{sec:retrieval-results}

Table~\ref{tab:retrieval-recall} measures whether any top-five passage contains
a gold answer or alias. The decomposition-path query policy substantially
improves recall on the three multi-hop datasets and leaves \popqa{} almost
unchanged. The largest gains occur on \twowiki{} and \musique{}, where bridge
entities make the later query more specific. As noted in
Section~\ref{sec:experiments}, the two paths also differ in title matching, so
these values should not be attributed to decomposition wording alone.

\begin{table}[t]
\caption{Answer-containing Hit@5 for original and decomposition-path queries.}
\label{tab:retrieval-recall}
\centering
\small
\setlength{\tabcolsep}{5pt}
\begin{tabular}{lrccc}
\toprule
Dataset & \(N\) & Original & Decomposition path & Difference \\
\midrule
\twowiki{} & 917 & 32.61 & \textbf{60.63} & +28.03 \\
\hotpot{} & 941 & 47.40 & \textbf{52.92} & +5.53 \\
\musique{} & 1,000 & 19.00 & \textbf{38.80} & +19.80 \\
\popqa{} & 1,000 & \textbf{71.20} & 71.10 & -0.10 \\
\bottomrule
\end{tabular}
\end{table}

The gate controls an accuracy--retrieval trade-off rather than eliminating
most retrieval calls. At the accuracy-first point, it is tied with always
retrieve to two decimal places while retrieving \(93.84\%\) of steps. Raising
the threshold to \(0.008461\) lowers retrieval to \(87.82\%\) at a cost of
\(0.26\) macro EM and \(0.24\) macro F1. More aggressive skipping causes
larger degradation (Table~\ref{tab:gate-sweep}). These are fixed-chain
replays: stored step branches are reselected and the chain is resynthesized,
while later resolved sub-questions remain fixed.

\begin{table}[t]
\caption{Accuracy--retrieval trade-off in the fixed-chain gate replay.}
\label{tab:gate-sweep}
\centering
\small
\setlength{\tabcolsep}{8pt}
\begin{tabular}{lcc}
\toprule
Policy / threshold & Macro EM/F1 & Step retrieval \\
\midrule
Never retrieve & 25.98/32.34 & 0.00\% \\
\(\tau=0.002305\) & \textbf{41.88}/\textbf{48.90} & 93.84\% \\
\(\tau=0.008461\) & 41.62/48.66 & 87.82\% \\
\(\tau=0.044170\) & 40.60/47.72 & 75.51\% \\
\(\tau=0.116500\) & 39.42/46.62 & 64.40\% \\
Always retrieve & \textbf{41.88}/\textbf{48.90} & 100.00\% \\
\bottomrule
\end{tabular}
\end{table}

\subsection{What does late routing add?}
\label{sec:routing-results}

The two candidates are complementary, but decomposition supplies more unique
correct answers on every dataset. Across the four \(1{,}000\)-question traces,
the direct path alone is correct in \(173\) cases, whereas only the
decomposition path is correct in \(874\) cases. Relative to always choosing
decomposition, the verifier recovers \(52\) direct-only cases but loses \(41\)
decomposition-only cases, for a net \(+11\) correct answers
(\(+0.28\) macro EM). The gain is concentrated on \hotpot{}; routing changes
EM by \(-0.4\), \(-0.2\), and \(-0.1\) on \twowiki{}, \musique{}, and
\popqa{}, respectively. We therefore view the verifier as a conservative
fallback, not the source of the main performance difference. Full overlap and
oracle counts appear in Appendix~\ref{app:routing}.

\section{Limitations}
\label{sec:limitations}

The present evidence is diagnostic rather than a complete confirmatory
evaluation. First, the operating threshold was selected on the same
development traces used for reporting; a disjoint calibration/evaluation split
is needed before drawing general conclusions about the selected point. Second,
all results use one backbone, greedy decoding, and one run on the first
\(1{,}000\) development examples. The \musique{} slice contains only two-hop
questions and understates the difficulty of the complete development set.
Third, external baselines are paper-reported rather than reproduced in the same
harness. Fourth, original and decomposition-path retrieval differ in title
matching, and original-question evidence union has not been isolated in a
separate ablation. The retrieval-recall gain therefore belongs to the complete
query policy, not uniquely to decomposition.

The threshold sweep is a fixed-chain replay rather than an independent live
rerun at every point. It preserves a controlled comparison but cannot capture
all downstream changes caused by an earlier retrieval decision. Finally,
\tailleak{} measures the shape of a model distribution, not factual
correctness. It cannot flag a false answer generated with extreme confidence,
and routing cannot repair questions for which both candidates are wrong.

\section{Conclusion}
\label{sec:conclusion}

We introduced \tailleak{}, a training-free confidence statistic that isolates
probability mass leaving a model's leading token set and accumulates it across
an answer. In the evaluated Llama-3-8B traces, it ranks direct-answer errors
better than representative entropy, likelihood, disagreement, margin, and
self-evaluation signals, and it serves as a useful sub-question retrieval gate
inside \method{}. The results also set a clear boundary: most decomposed steps
still benefit from retrieval, late routing adds only a small correction, and
confident hallucinations remain. Broader backbones, held-out calibration, and
controlled retrieval ablations are the next tests required to establish how
far the signal generalizes.
\label{sec:main-end}


\subsection*{AI use statement}
Generative AI tools were used to help organize the manuscript, draft and edit
prose, translate between Chinese and English, review code, format \LaTeX{}, and
assist with interpreting already-computed experimental results. The reported
numbers were checked against saved experimental artifacts, and all AI-assisted
text and code were reviewed by the authors. The authors take responsibility for
the final content, including text and artifacts produced with AI assistance.

\subsection*{Ethics statement}
This work evaluates language models on existing public question-answering
benchmarks and Wikipedia passages. It introduces no new human-subject data.
Like other retrieval-augmented generation systems, the method can reproduce
biases or factual errors present in the model and retrieval corpus; confidence
scores should not be treated as guarantees of correctness.

\subsection*{Reproducibility statement}
Section~\ref{sec:experiments} specifies the model, data slices, corpus,
retrieval configuration, training setup, and metrics. Appendix~\ref{app:details}
provides implementation details, Appendix~\ref{app:additional-results} gives
additional retrieval analyses, and Appendices~\ref{app:confidence}
and~\ref{app:routing} report confidence and routing diagnostics. The
supplementary package is intended to include the scripts and saved traces
needed to reproduce the reported tables.

\bibliography{../shared/references}
\bibliographystyle{iclr2027_conference}

\appendix
\section{Implementation Details}
\label{app:details}

\subsection{Training and inference configuration}

\begin{table}[h]
\caption{Configuration used for the reported \method{} runs.}
\label{tab:implementation}
\centering
\small
\setlength{\tabcolsep}{6pt}
\begin{tabular}{ll}
\toprule
Setting & Value \\
\midrule
Backbone & Llama-3-8B-Instruct \\
Decomposer data & 12,000 examples; 4,000 per source \\
LoRA & rank 16, scale 32, dropout 0.05 \\
Target modules & all linear modules \\
Training & 3 epochs, bf16, effective batch 16 \\
Optimization & \(10^{-4}\), cosine decay, 10\% warmup \\
Maximum context & 1,024 tokens \\
Decoding & greedy; at most 48 answer tokens \\
Retrieval & BM25, top 8, Wikipedia-21M \\
Tail set & top \(k=5\) tokens \\
Step gate & \(\tau=0.002305\) \\
Direct verifier & \(\tau_d=0.0005\) \\
Post-decomposition verifier & \(\tau_p=0.0158\) \\
Type-aware margin & 2.5 \\
\bottomrule
\end{tabular}
\end{table}

\subsection{Prompts}

The answer probe uses the same answer-only instruction for direct and
sub-question generation:
\begin{quote}
\small
\emph{Answer the question with the answer phrase only. No explanation, no
preamble, no full sentence---just the answer.}
\end{quote}
For a decomposed step, the original question is prepended when it differs from
the resolved sub-question. Retrieved context, when available, appears between
the original question and sub-question.

The decomposer instruction asks for an ordered list of atomic sub-questions,
one per line, with the format \(\texttt{\#<index>:}\). Dependencies must use a
\(\texttt{\#N}\) placeholder. The final synthesizer receives the original
question and all sub-question/answer pairs. It returns only a short answer span
and contains explicit rules for comparison, yes/no, and year questions. The
supplementary source package contains the exact strings and chat-template
construction.

\subsection{Final routing rule}

The final rule is deterministic. It first handles exact candidate agreement,
empty or malformed outputs, unresolved placeholders, and obvious yes/no or
comparison-type mismatches. If these rules do not decide the route, it chooses
direct only when
\[
 \tl(a^{\mathrm{dir}})<\tau_d
 \quad\text{and}\quad
 \max_i\tl(a_i^{\mathrm{post}})>\tau_p
\]
or the decomposition is structurally broken; otherwise it chooses
decomposition. A final type-aware check is applied only to undecided
decomposition defaults. It scores answer-type compatibility, whether the final
answer repeats an intermediate answer, whether a step answer appears as an
entity fragment inside a later query, and whether a candidate simply repeats
part of the question. It switches to direct only when the direct score exceeds
the decomposition score by at least 2.5.

\section{Additional Retrieval Results}
\label{app:additional-results}

\subsection{Gold and untuned decompositions on MuSiQue}

The first 1,000 \musique{} development questions are all two-hop, allowing
comparison with the official decomposition. Table~\ref{tab:gold-decomp}
shows that \method{} improves over the original query but remains below a gold
chain. Replacing the bridge entity with the gold answer closes most of the gap,
indicating that bridge-answer drift is more important than sub-question wording
in this slice.

\begin{table}[h]
\caption{Answer-containing Hit@5 on the MuSiQue slice.}
\label{tab:gold-decomp}
\centering
\small
\setlength{\tabcolsep}{7pt}
\begin{tabular}{lc}
\toprule
Query condition & Hit@5 \\
\midrule
Original question & 19.0 \\
\method{} decomposition & 38.8 \\
Gold decomposition + gold bridge & \textbf{55.8} \\
Gold decomposition, unresolved \(\texttt{\#1}\) & 19.8 \\
\method{} wording + gold bridge & 54.1 \\
Untuned base wording + gold bridge & 52.9 \\
\bottomrule
\end{tabular}
\end{table}

\subsection{Title matching}

Original questions are long and constraint-rich, while decomposed queries are
short and entity-centered. Table~\ref{tab:title-policy} compares title matching
for both query types. These are live end-to-end reruns.

\begin{table}[h]
\caption{EM/F1 under three BM25 title policies.}
\label{tab:title-policy}
\centering
\small
\setlength{\tabcolsep}{7pt}
\begin{tabular}{lccc}
\toprule
Dataset & Titles off & Titles on & Path-aware \\
\midrule
\popqa{} & 51.9/54.0 & 52.4/54.6 & \textbf{53.4/55.7} \\
\hotpot{} & 37.3/48.7 & 37.1/48.4 & \textbf{37.9/49.3} \\
\musique{} & 23.3/31.6 & 24.9/\textbf{34.5} & \textbf{26.2}/34.4 \\
\twowiki{} & 44.5/50.4 & 49.1/55.5 & \textbf{49.6/56.3} \\
\bottomrule
\end{tabular}
\end{table}

\section{Additional Confidence Results}
\label{app:confidence}

\subsection{Signal definitions}

All scores are oriented so larger values indicate greater uncertainty.
\begin{itemize}
    \item \textbf{Maximum token entropy}: the largest next-token entropy in
    the generated answer.
    \item \textbf{Sequence NLL}: the summed negative log-likelihood of the
    generated tokens.
    \item \textbf{TARG margin}: the transformed average gap between the two
    largest logits over the first eight answer tokens.
    \item \textbf{Sample disagreement}: the fraction of ten sampled answers
    that differ from the greedy answer.
    \item \textbf{P(True)}: one minus the model's estimated probability that
    its proposed answer is true.
    \item \textbf{\tailleak{}}: the sum of negative log top-five mass in
    Equation~\ref{eq:tail-leakage}.
\end{itemize}
The first three and \tailleak{} use the same greedy pass. Sample disagreement
and P(True) require additional generations.

\subsection{Entropy components}

Table~\ref{tab:entropy-components} evaluates the three terms of
Equation~\ref{eq:entropy-decomposition}. The boundary-crossing component
carries most of the useful signal; the arrangement within the tail is close to
chance. These component values come from the saved diagnostic computation and
should be recomputed from complete logits before being elevated to a primary
claim.

\begin{table}[h]
\caption{Macro AUROC of entropy components.}
\label{tab:entropy-components}
\centering
\small
\begin{tabular}{lc}
\toprule
Component & Macro AUROC \\
\midrule
Leading set versus tail & 0.742 \\
Competition within top five & 0.712 \\
Arrangement within tail & 0.502 \\
Full token entropy & 0.726 \\
Answer-level \tailleak{} & \textbf{0.749} \\
\bottomrule
\end{tabular}
\end{table}

\subsection{End-to-end signal replacement}

\begin{table}[h]
\caption{EM/F1 with entropy and Tail Leakage.}
\label{tab:entropy-vs-tail}
\centering
\scriptsize
\setlength{\tabcolsep}{3.5pt}
\begin{tabular}{lccccc}
\toprule
Signal & \hotpot{} & \twowiki{} & \musique{} & \popqa{} & Macro \\
\midrule
Max entropy & 35.50/46.11 & \textbf{50.00}/\textbf{56.58} & 24.70/33.06 & 51.90/54.71 & 40.52/47.61 \\
\tailleak{} & \textbf{37.80}/\textbf{49.03} & 49.30/56.22 & \textbf{26.50}/\textbf{34.63} & \textbf{53.90}/\textbf{55.73} & \textbf{41.88}/\textbf{48.90} \\
\bottomrule
\end{tabular}
\end{table}

\subsection{Selective answering}

Table~\ref{tab:selective} ranks direct answers by \tailleak{} and retains only
the most confident fraction. Accuracy increases as coverage decreases, but the
hard multi-hop datasets remain unreliable even at low coverage.

\begin{table}[h]
\caption{Direct-path EM by confidence coverage.}
\label{tab:selective}
\centering
\small
\setlength{\tabcolsep}{6pt}
\begin{tabular}{lccccc}
\toprule
Coverage & \hotpot{} & \twowiki{} & \musique{} & \popqa{} & Pooled \\
\midrule
25\% & 57.6 & 34.0 & 22.4 & 79.2 & 52.8 \\
50\% & 45.4 & 26.8 & 15.2 & 57.0 & 38.9 \\
75\% & 38.0 & 22.1 & 12.5 & 46.3 & 30.2 \\
100\% & 30.3 & 18.0 & 10.2 & 37.8 & 24.1 \\
\bottomrule
\end{tabular}
\end{table}

\section{Routing and Failure Analysis}
\label{app:routing}

\subsection{Candidate overlap}

\begin{table}[h]
\caption{Correctness overlap before final routing.}
\label{tab:path-overlap}
\centering
\small
\setlength{\tabcolsep}{5pt}
\begin{tabular}{lrrrr}
\toprule
Dataset & Both correct & Direct only & Decomp. only & Neither \\
\midrule
\hotpot{} & 225 & 78 & 135 & 562 \\
\twowiki{} & 144 & 36 & 353 & 467 \\
\musique{} & 75 & 27 & 192 & 706 \\
\popqa{} & 346 & 32 & 194 & 428 \\
\bottomrule
\end{tabular}
\end{table}

\subsection{Verifier effect}

\begin{table}[h]
\caption{Verifier changes relative to always choosing decomposition.}
\label{tab:verifier-effect}
\centering
\small
\setlength{\tabcolsep}{5pt}
\begin{tabular}{lrrr}
\toprule
Dataset & Direct-only recovered & Decomp.-only lost & Net EM \\
\midrule
\hotpot{} & 39/78 & 21/135 & +1.80 \\
\twowiki{} & 3/36 & 7/353 & -0.40 \\
\musique{} & 9/27 & 11/192 & -0.20 \\
\popqa{} & 1/32 & 2/194 & -0.10 \\
\midrule
Pooled / macro & 52/173 & 41/874 & +0.28 \\
\bottomrule
\end{tabular}
\end{table}

\subsection{Decomposition-path failures}

Each incorrect multi-hop decomposition candidate is assigned one primary
cause. Evidence is counted as present when the required Wikipedia title
appears in the top ten results from either the original or resolved
sub-question queries.

\begin{table}[h]
\caption{Primary causes of incorrect decomposition answers.}
\label{tab:failure-taxonomy}
\centering
\scriptsize
\setlength{\tabcolsep}{3.5pt}
\begin{tabular}{lccc}
\toprule
Cause & \hotpot{} (\(N=640\)) & \twowiki{} (\(N=503\)) & \musique{} (\(N=733\)) \\
\midrule
Missing required evidence & 307 (48.0\%) & 311 (61.8\%) & 513 (70.0\%) \\
Wrong step answer & 131 (20.5\%) & 98 (19.5\%) & 73 (10.0\%) \\
Wrong final synthesis & 164 (25.6\%) & 92 (18.3\%) & 138 (18.8\%) \\
Lost original constraint & 38 (5.9\%) & 2 (0.4\%) & 9 (1.2\%) \\
\bottomrule
\end{tabular}
\end{table}

Missing evidence is the largest category on all three multi-hop datasets.
Routing cannot repair many of these cases because the direct candidate is also
wrong.

\subsection{Confidently wrong cases}

\begin{table}[h]
\caption{Tail Leakage when both candidates agree in correctness.}
\label{tab:quadrant-confidence}
\centering
\small
\setlength{\tabcolsep}{4.5pt}
\begin{tabular}{lcccc}
\toprule
Dataset & Both-correct mean & Both-wrong mean & AUROC & Confidently wrong \\
\midrule
\hotpot{} & 0.01260 & 0.12927 & 0.777 & 15.5\% \\
\twowiki{} & 0.03414 & 0.20215 & 0.751 & 20.1\% \\
\musique{} & 0.04945 & 0.25025 & 0.766 & 12.7\% \\
\popqa{} & 0.00974 & 0.10282 & 0.846 & 4.9\% \\
Pooled & 0.01877 & 0.17926 & 0.817 & 9.3\% \\
\bottomrule
\end{tabular}
\end{table}

Both-wrong examples have higher scores on average, but the groups overlap.
This is why \tailleak{} is used as a population-level ranking and gating signal
rather than a per-question correctness test.


\end{document}
